\documentclass[QCM, noCustomPackages, basic]{UPSTI_Document}


\UPSTIloadProfile{BUT_GEII}

\newcommand{\sequence}{2}
\newcommand{\UPSTIsousTitreEnTete}{Sequence \sequence~: Boucles et fonctions}


\graphicspath{{images/}}



\discipline{Informatique}
\newcommand{\UPSTIidClasse}{1}
\newcommand{\UPSTItitreEnTete}{Informatique}


%\newcommand{\UPSTIsource}{Certaines images proviennent du site freepics.com}

\usepackage{listings}
\usepackage{xcolor}
\usepackage{amssymb} % pour \Box et \CheckedBox
\usepackage{makecell} % permet \makecell{...} avec des \\ dans une cellule
\renewcommand\arraystretch{1.2}

\lstdefinelanguage{bash}{
  alsoletter={-},           % autoriser "-" dans les mots-clés
  morekeywords={
    ls, cd, mv, cp, rm, mkdir, rmdir, echo, cat, less, head, tail,
    grep, find, chmod, chown, ps, kill, top, man, which, alias, export,
    ssh, scp, rsync, tar, gzip, gunzip, ping, curl, wget, zip, unzip, code
  },
  sensitive=true,           % sensible à la casse
  morecomment=[l]{\#},      % commentaires avec #
  morestring=[b]",          % chaînes "..."
  morestring=[b]'           % chaînes '...'
}

\definecolor{colString}{rgb}{0.6,0.1,0.1}
\definecolor{lightblack}{rgb}{0.3,.3,.3}

\newcommand{\prompt}[1]{\textcolor{gray}{#1}}

% \lstdefinestyle{console}{
%   language=bash,
%   basicstyle=\ttfamily,
%   columns=fullflexible,
%   keepspaces=true,
%   escapeinside=||, % <- permet d'injecter du LaTeX entre |...|
% }

\lstdefinestyle{console}{
    basicstyle      = \footnotesize\ttfamily,% ecriture standard
    extendedchars   = true,% permet d'avoir des accents dans le code
    tabsize         = 2,%
    showspaces      = false,%
    showstringspaces = false,%
    numbers=none,%
    breaklines=true,%
    breakautoindent=true,%
    escapeinside=|| % <- permet d'injecter du LaTeX entre |...|
}




\def\photodir{photos}

%----------------------------------------------------------------------
% \trou{largeur} : cadre pointillé inséré DANS un listing via
% escapeinside={(*}{*)}, à l'endroit exact où l'étudiant doit écrire.
% Exemple d'usage dans un fichier .c :  float duree_h = (*\trou{4cm}*);
%----------------------------------------------------------------------
\newcommand{\trou}[1]{%
    \tikz[baseline=(trou.base)]{%
        \node[draw=gray!70, dashed, rounded corners=2pt, inner sep=2pt,
              minimum width=#1, minimum height=2.4em, fill=white,
              anchor=base] (trou) {\strut};%
    }%
}


%----------------------------------------------------------------------
% Numérotation automatique des parties et sous-parties :
%   \section{Sujet}     ->  « Partie A — Sujet »
%   \subsection{Thème}  ->  « A1 — Thème »  (remis à zéro à chaque partie)
%----------------------------------------------------------------------
\renewcommand{\thesection}{\Alph{section}}
\renewcommand{\thesubsection}{\thesection\arabic{subsection}}
\makeatletter
\def\DS@fmt@section{Partie~\thesection~--- }
\def\DS@fmt@subsection{\thesubsection~--- }
\renewcommand{\@seccntformat}[1]{%
    \ifcsname DS@fmt@#1\endcsname\csname DS@fmt@#1\endcsname
    \else\csname the#1\endcsname\quad\fi}
\makeatother

% Masque le libellé "Question n" (grilles où chaque ligne est une question AMC ;
% le numéro affiché est géré par le compteur UPSTIopenQNum du package UPSTI_AMC)
\newcommand{\sansNumeroQuestion}{\def\AMCbeginQuestion##1##2{}}
% Affiche le libellé de la prochaine question sans la consommer
\newcommand{\libelleProchaineQuestion}{\textbf{Question \the\numexpr\value{UPSTIopenQNum}+1\relax}}
% Après une grille de questions AMC sans libellé, la grille compte pour une question affichée
\newcommand{\grilleCommeUneQuestion}{\stepcounter{UPSTIopenQNum}}

\documentVersion{E}
\newcommand{\UPSTInumeroVersion}{0}
\newcommand{\UPSTImessage}{DS1 - Borne de recharge}
\newcommand{\UPSTItitre}{Variables, conditions, boucles et fonctions}

\newcommand{\sujetA}{
\exemplaire{3}{%
\noindent\champnom{\fbox{
    \begin{minipage}{.9\linewidth}
        Nom, prénom et groupe :

        \vspace*{.5cm}\dotfill
        \vspace*{1mm}
    \end{minipage}
}}

\renewcommand{\UPSTImessage}{Sujet A}
\renewcommand{\UPSTInumeroVersion}{A}
\UPSTIbuildPage
\setcounter{figure}{0}
\setcounter{table}{0}

%======================================================================
% CONTEXTE GÉNÉRAL — cahier des charges du système étudié
%======================================================================
\section*{Contexte : la borne de recharge \textsc{EV-BR150}}

Vous intégrez le bureau d'études d'un équipementier qui développe la borne de recharge
pour véhicules électriques \textsc{EV-BR150}, destinée aux parkings d'entreprise. Votre
mission porte sur le \textbf{logiciel embarqué} de supervision de la borne : il mesure,
décide, facture et rend compte. Le cahier des charges qui vous est remis est le suivant.

\medskip
\noindent\textbf{Architecture et puissance.} La borne délivre, par point de charge, une
puissance qui va de 3,7 kW à \textbf{150 kW} au maximum. Trois modes d'exploitation sont définis selon la
puissance demandée par le véhicule : mode \emph{lent} jusqu'à 22 kW inclus, mode
\emph{rapide} au-delà de 22 kW et jusqu'à 50 kW inclus, mode \emph{ultra-rapide} au-delà de
50 kW. Toute demande supérieure à 150 kW est hors du domaine de la borne.

\medskip
\noindent\textbf{Grandeurs électriques et sécurité.} Le raccordement est triphasé
400 V ; la tension réseau mesurée doit rester comprise entre \textbf{360 V et 440 V}
(bornes incluses). Le courant demandé ne doit pas dépasser \textbf{63 A}. La température du
module de puissance doit rester \textbf{strictement inférieure à 65 °C}. La charge ne
s'engage que si le câble est \textbf{verrouillé} côté borne. Dès qu'une de ces conditions
n'est plus respectée, la charge est interrompue.

\medskip
\noindent\textbf{Tarification.} L'énergie est facturée 0,22 €/kWh en mode lent et rapide,
et 0,30 €/kWh au-delà de 50 kW. Un tarif heures creuses de 0,15 €/kWh s'applique aux
sessions démarrées entre 22 h et 5 h inclus. Chaque point de charge utilisé ajoute
\textbf{0,50 € de frais de service} par session. Le stationnement au-delà de 120 minutes
peut être majoré.

\medskip
\noindent\textbf{Exploitation.} Les durées sont relevées en minutes par l'automate, alors
que les tarifs sont exprimés au kWh : les conversions sont à la charge du logiciel. En fin
de journée, l'exploitant attend un bilan : nombre de sessions valides, énergie délivrée,
recette. Une session est déclarée \textbf{anormale} (donc ni comptée, ni facturée) si sa
puissance dépasse 150 kW ou si sa durée est nulle ou négative.

\medskip
\noindent\textbf{Déroulement du sujet.} Les parties A à E suivent les étapes de conception de
ce logiciel : on modélise d'abord les \textbf{grandeurs} (partie A), puis les \textbf{règles
de sécurité et de mode} (partie B), puis le \textbf{traitement d'une série de sessions}
(partie C), puis la \textbf{modularisation en fonctions} (partie D), et enfin
l'\textbf{intégration} de l'ensemble (partie E). Les cadres pointillés \trou{1.6cm}
placés dans les programmes sont à compléter directement sur le sujet.

\medskip
\noindent\textbf{Reprise du code existant.} Une première version de ce logiciel a été écrite
par un développeur débutant qui vient de quitter l'entreprise. Vous êtes recruté pour
\textbf{reprendre et fiabiliser} son travail. Tous les extraits de code du sujet proviennent
de cette version et peuvent contenir des erreurs : \textbf{ne supposez jamais qu'un extrait est
correct}, déterminez ce que le programme fait \emph{réellement}, et non ce qu'il était censé
faire.

\pagebreak

\section{Modéliser les grandeurs de la session de charge}

\subsection{Déclarer des variables}
\restituegroupe{var-types}
\restituegroupe{var-qcm-declaration}
\subsection{Calcul des heures entamées}
\restituegroupe{var-qcm-conversion}
\subsection{Le dépassement}
\restituegroupe{var-qcm-priorite}
\needspace{20\baselineskip}% garde le titre avec son code
\subsection{Trace d'un calcul de coût}
\restituegroupe{var-trace}
%\restituegroupe{var-completion-decl}

\pagebreak

\section{Décider : sécurité et mode de charge}

La borne doit décider : autoriser ou refuser la charge, surveiller la sécurité.


\subsection{Lire une condition}
\restituegroupe{cond-lecture-if}
\subsection{Autoriser la charge}
\restituegroupe{cond-completion}
\subsection{Surveiller la sécurité}
\restituegroupe{cond-correction}

\pagebreak

\section{Répéter : traiter une série de sessions}

La borne répète sans cesse des traitements : attentes, cumuls, simulations.

\subsection{Attente du verrouillage}
\restituegroupe{boucle-qcm-borne}
\needspace{24\baselineskip}% garde le titre avec son code
\subsection{Cumul de l'énergie}
\restituegroupe{boucle-correction}
\needspace{22\baselineskip}% garde le titre avec son code
\subsection{Recharge à puissance dégressive}
\restituegroupe{boucle-trace}
\needspace{38\baselineskip}% garde le titre avec son code
\subsection{Montée en charge par paliers}
\restituegroupe{boucle-completion}
%\restituegroupe{boucle-conception}

\pagebreak

\section{Modulariser : des fonctions pour la borne}

Le logiciel grossit : on isole ses traitements dans des fonctions réutilisables.

\needspace{34\baselineskip}% garde le titre avec son code
\subsection{Appeler une fonction existante}
\restituegroupe{fonc-lecture-appel}
\needspace{30\baselineskip}% garde le titre avec son code
\subsection{Portée d'une variable}
\restituegroupe{fonc-qcm-portee}
% \restituegroupe{fonc-correction-portee}
\needspace{40\baselineskip}% garde le titre avec son code
\subsection{Écrire une fonction de décision}
\restituegroupe{fonc-ecrire-mode}

\pagebreak

\section{Intégrer : le bilan de journée}
\restituegroupe{synthese-utilisation}
% \restituegroupe{synthese-evolution}
% \restituegroupe{synthese-critique}

% \pagebreak

% \section{Questions supplémentaires}

% Ces deux questions, plus délicates, portent sur des subtilités du langage C (représentation
% des nombres à virgule, évaluation des conditions composées).

% \restituegroupe{var-qcm-flottant}
% \restituegroupe{cond-qcm-courtcircuit}

\AMCcleardoublepage
}
}


\begin{document}

\vspace{1em}

\begin{center}
    \textit{Aucun document autorisé, sans calculatrice, durée : 55 min}\\
    \textit{Si vous constatez une erreur dans le document, indiquez les suppositions que vous faites pour répondre à la question.}\\
    \textit{Chaque partie porte sur ses propres données et extraits de code, indépendamment des autres parties.}
\end{center}

\setdefaultgroupmode{withoutreplacement}
% ==================== PARTIE A — VARIABLES ====================
% Étape de conception : modéliser les grandeurs de base d'une session de charge
% (puissance, durée, énergie, tarif, montant) et choisir les bons types.
%
% Progression :
%   A1  trace de valeurs (tableau à compléter)
%   A2  QCM — division entière et modulo (conversion min -> h)
%   A3  QCM — priorité des opérateurs (division exacte) + expression générale
%   A4a compléter les déclarations (types imposés)
%   A4b compléter les 3 expressions de calcul
%   (l'ancien QCM flottant est déplacé en partie F)

\element{var-types}{
    \noindent\libelleProchaineQuestion\\
    Pour chacune des grandeurs ci-dessous manipulées par la borne, cocher \textbf{une seule case} :
    le type le plus adapté pour la variable qui la représente.
    (\texttt{char} et \texttt{string} ne sont proposés ici qu'à titre d'options de reconnaissance.)

    \begin{center}
    {\sansNumeroQuestion\renewcommand{\arraystretch}{1.5}
    \begin{tabular}{|p{5.6cm}|p{10.6cm}|}
        \hline
        \textbf{Grandeur} & \textbf{Type C le plus adapté} \\ \hline
        Nombre de badges utilisateurs enregistrés dans la borne &
        \begin{question}{var-types-nbsessions}\begin{reponseshoriz}[o] \bonne{\texttt{int}} \mauvaise{\texttt{float}} \mauvaise{\texttt{char}} \mauvaise{\texttt{string}} \end{reponseshoriz}\end{question} \\ \hline
        Tension réseau mesurée (ex. 398,6 V) &
        \begin{question}{var-types-tension}\begin{reponseshoriz}[o] \mauvaise{\texttt{int}} \bonne{\texttt{float}} \mauvaise{\texttt{char}} \mauvaise{\texttt{string}} \end{reponseshoriz}\end{question} \\ \hline
        Nombre de coupures réseau détectées depuis la mise en service &
        \begin{question}{var-types-point}\begin{reponseshoriz}[o] \bonne{\texttt{int}} \mauvaise{\texttt{float}} \mauvaise{\texttt{char}} \mauvaise{\texttt{string}} \end{reponseshoriz}\end{question} \\ \hline
        Rendement du convertisseur de puissance (ex. 0,94) &
        \begin{question}{var-types-tarif}\begin{reponseshoriz}[o] \mauvaise{\texttt{int}} \bonne{\texttt{float}} \mauvaise{\texttt{char}} \mauvaise{\texttt{string}} \end{reponseshoriz}\end{question} \\ \hline
        Un caractère représentant la catégorie d'un utilisateur (ex. 'A', 'B', 'C') &
        \begin{question}{var-types-duree}\begin{reponseshoriz}[o] \mauvaise{\texttt{int}} \mauvaise{\texttt{float}} \bonne{\texttt{char}} \mauvaise{\texttt{string}} \end{reponseshoriz}\end{question} \\ \hline
        Index du compteur d'énergie depuis l'installation (ex. 18\,452,7 kWh) &
        \begin{question}{var-types-energie}\begin{reponseshoriz}[o] \mauvaise{\texttt{int}} \bonne{\texttt{float}} \mauvaise{\texttt{char}} \mauvaise{\texttt{string}} \end{reponseshoriz}\end{question} \\ \hline
    \end{tabular}}
    \end{center}
    \grilleCommeUneQuestion
}

\element{var-qcm-declaration}{
    \begin{question}{var-qcm-declaration}
        Quelle ligne déclare correctement une variable entière \lstinline[language=C]|nb_sessions|
        initialisée à zéro ?
        \begin{multicols}{3}
            \begin{reponses}
                \bonne{\lstinline[language=C]|int nb_sessions = 0;|}
                \mauvaise{\lstinline[language=C]|nb_sessions int = 0;|}
                \mauvaise{\lstinline[language=C]|int nb_sessions == 0;|}
                \mauvaise{\lstinline[language=C]|int 0 = nb_sessions;|}
                \mauvaise{\lstinline[language=C]|(nb_sessions == 0);|}
                \mauvaise{\lstinline[language=C]|int nb_sessions == 0;|}
            \end{reponses}
        \end{multicols}
    \end{question}
}

\element{var-trace}{
    \begin{question}{var-trace-1}
        Avant d'écrire le logiciel, on vérifie la façon dont la borne manipule ses grandeurs.
        Le fragment ci-dessous traite une session de 7 kW pendant 150 minutes,
        facturée 0,22 €/kWh.
        
        Compléter le tableau en donnant la valeur de la variable affectée par chacune des lignes
        repérées \texttt{L1} à \texttt{L5}, \textbf{en respectant le type de la variable}
        (une valeur entière s'écrit sans décimale).
        \evaluationProf[2]

        \begin{minipage}[t]{.56\textwidth}
            \lstinputlisting[language=c]{sources/codes/a1-trace-grandeurs.c}
        \end{minipage}\hfill
        \begin{minipage}[t]{.42\textwidth}~\\
            {\renewcommand{\arraystretch}{2}
            \begin{tabular}{|c|l|p{2.6cm}|}
                \hline
                \textbf{Ligne} & \textbf{Variable} & \textbf{Valeur} \\ \hline
                L1 & \lstinline[language=C]|duree_h| & \\ \hline
                L2 & \lstinline[language=C]|energie_kwh| & \\ \hline
                L3 & \lstinline[language=C]|energie_kwh_ent|  & \\ \hline
                L4 & \lstinline[language=C]|energie_kwh| & \\ \hline
                L5 & \lstinline[language=C]|cout_euros|  & \\ \hline
            \end{tabular}}

            \medskip
        \end{minipage}
    \end{question}
}

\element{var-qcm-nomVariables}{
    Le développeur a choisi des noms de variables qui ne sont pas très explicites.
    Il souhaite les renommer pour mieux refléter leur rôle dans le calcul du coût
    d'une session de charge.
    
    \begin{center}
        \begin{minipage}{.74\textwidth}
            \lstinputlisting[language=c]{sources/codes/a2-qcm-nomVariables.c}
        \end{minipage}
    \end{center}
    
    \begin{question}{var-qcm-nomVariables}
        Quel nom de variable est le plus approprié pour remplacer \lstinline[language=C]|a| ?
        \begin{multicols}{4}
            \begin{reponses}
                \bonne{\lstinline[language=C]|duree_min|}
                \mauvaise{\lstinline[language=C]|heures_entamees|}
                \mauvaise{\lstinline[language=C]|reste_min|}
                \mauvaise{\lstinline[language=C]|minutes_restant|}
            \end{reponses}
        \end{multicols}
    \end{question}
}

\element{var-qcm-conversion}{
    Pour le ticket de stationnement, la borne convertit la durée (min) en
    \emph{heures entamées} et en minutes restantes.

    \begin{center}
        \begin{minipage}{.74\textwidth}
            \lstinputlisting[language=c]{sources/codes/a2-qcm-conversion.c}
        \end{minipage}
    \end{center}

    \begin{question}{var-lecture-printf}
        Pour \lstinline[language=C]|heures_entamees|
        valant \texttt{2} et \lstinline[language=C]|reste_min| valant \texttt{7}, qu'affiche la dernière
        ligne (l'appel à \lstinline[language=C]|printf|) du programme ci-dessous ?
        \begin{multicols}{2}
            \begin{reponses}
                \bonne{\texttt{2 7}}
                \mauvaise{\texttt{7 2}}
                \mauvaise{\texttt{27}}
                \mauvaise{\texttt{722}}
                \mauvaise{\texttt{\%d \%d}}
                \mauvaise{\texttt{heures\_entamees reste\_min}}
            \end{reponses}
        \end{multicols}
    \end{question}
    


    \begin{question}{var-qcm-conversion}
        Qu'affiche ce fragment de code complet (avec \lstinline[language=C]|duree_min| égal à 185) ?
        \begin{multicols}{6}
            \begin{reponses}
                \bonne{\texttt{4 5}}
                \mauvaise{\texttt{3 5}}
                \mauvaise{\texttt{4 65}}
                \mauvaise{\texttt{3 65}}
                \mauvaise{\texttt{4 55}}
                \mauvaise{\texttt{3 55}}
            \end{reponses}
        \end{multicols}
    \end{question}
}

\element{var-qcm-priorite}{
    Le stationnement est majoré de 3~\euro{} par heure entamée \emph{au-delà} de
    120 minutes, en plus de 50 centimes de frais de service. 
    
    Un développeur inexpérimenté écrit :
    
    \begin{center}
        \begin{minipage}{.74\textwidth}
            \lstinputlisting[language=c]{sources/codes/a3-qcm-priorite.c}
        \end{minipage}
    \end{center}
    
    \begin{question}{var-qcm-priorite-180}
        Qu'affiche ce fragment de code si on entre une durée de 180 minutes ?
        \begin{multicols}{6}
            \begin{reponses}
                \bonne{350 cts}
                \mauvaise{50 cts}
                \mauvaise{200 cts}
                \mauvaise{950 cts}
                \mauvaise{0 cts}
                \mauvaise{300 cts}
            \end{reponses}
        \end{multicols}
    \end{question}

    \begin{question}{var-qcm-priorite-150}
        Qu'affiche ce fragment de code si on entre une durée de 150 minutes ?
        \begin{multicols}{6}
            \begin{reponses}
                \mauvaise{350 cts}
                \bonne{50 cts}
                \mauvaise{200 cts}
                \mauvaise{950 cts}
                \mauvaise{0 cts}
                \mauvaise{300 cts}
            \end{reponses}
        \end{multicols}
    \end{question}

    \begin{question}{var-qcm-priorite-corrige}
        
        Proposer une correction du programme pour qu'il calcule correctement le montant à facturer.

        \evaluationProfGrille{2}{2}
    \end{question}
}

\element{var-completion-decl}{
    On facture maintenant une session réelle. La durée est saisie \textbf{en minutes},
    le tarif est exprimé en €/kWh, et tout montant comprend 0,50 € de frais de service.
    Les types sont imposés : \lstinline[language=C]|puissance_kw| est un nombre à virgule
    (\lstinline[language=C]|float|), \lstinline[language=C]|duree_min| un nombre entier
    (\lstinline[language=C]|int|), \lstinline[language=C]|tarif_kWh| un nombre à virgule
    (\lstinline[language=C]|float|).

    \textbf{Rappel :} énergie (kWh) $=$ puissance (kW) $\times$ durée (h).

    \lstinputlisting[language=c,escapeinside={(*}{*)}]{sources/codes/a4-cout-session-a-completer.c}

    \begin{question}{var-completion-decl}

        \evaluationProf[1]

        Compléter \textbf{directement dans les trois premiers cadres pointillés} la déclaration
        des variables \lstinline[language=C]|puissance_kw|, \lstinline[language=C]|duree_min|
        et \lstinline[language=C]|tarif_kWh|.

    \end{question}
    \begin{question}{var-completion-calcul}
        \evaluationProf[3]

        Dans le programme de la question précédente, compléter les \textbf{trois expressions
        de calcul} manquantes (cadres pointillés de \lstinline[language=C]|duree_h|,
        \lstinline[language=C]|energie_kwh| et \lstinline[language=C]|montant_euros|).
        Attention au type du résultat de chaque calcul intermédiaire.

        Utiliser, si besoin, l'espace ci-dessous pour justifier un choix de type ou de conversion.

    \end{question}
}

% ==================== PARTIE B — CONDITIONS ====================
% Étape de conception : règles de sécurité et choix du mode de charge.
%
% Progression :
%   B0  QCM — lire un if/else simple (sans piège)
%   B1  compléter UNE condition composée  (trou dans le code, escapeinside)
%   B2  repérer (numéroter) et corriger des erreurs
%   (QCM court-circuit déplacé en partie F ; conception supprimée, doublon de D)

\element{cond-lecture-if}{
  La borne surveille la température de son module de puissance avec le code suivant.

  \begin{question}{cond-lecture-if}
    \begin{minipage}{.45\textwidth}
      \lstinputlisting[language=c]{sources/codes/b0-lecture-if.c}
    \end{minipage}\hfill
    \begin{minipage}{.5\textwidth}
      Qu'affiche ce fragment de code ?
      \begin{reponses}
        \bonne{\texttt{Module OK}}
        \mauvaise{\texttt{Surchauffe : charge coupee}}
        \mauvaise{\texttt{Module OK} puis, sur la ligne suivante, \texttt{Surchauffe : charge coupee}}
        \mauvaise{Rien : aucune des deux branches n'est exécutée}
      \end{reponses}
    \end{minipage}
  \end{question}
}

\element{cond-completion}{
  \begin{question}{cond-completion}
    La borne n'autorise une charge que si \textbf{toutes} les conditions d'exploitation
    sont réunies :
  \begin{itemize}\setlength{\itemsep}{0pt}
      \item tension réseau comprise entre 360 V et 440 V, \textbf{bornes incluses} ;
      \item courant demandé \textbf{au plus} égal à 63 A ;
      \item température du module de puissance \textbf{strictement} inférieure à 65 °C.
    \end{itemize}

    Écrire la condition manquante \textbf{directement dans le cadre pointillé}.

    \lstinputlisting[language=c,escapeinside={(*}{*)}]{sources/codes/b2-autorisation-a-completer.c}

    Justifier brièvement le choix des opérateurs aux bornes de la plage de tension.

    \evaluationProfGrille{2}{3}
  \end{question}
}

\element{cond-correction}{
  \begin{question}{cond-correction}
    Un technicien a écrit le module de surveillance ci-dessous. Compilé sans option
    d'avertissement, il affiche parfois « Charge interrompue » alors qu'aucun défaut
    n'a été signalé.

    \lstinputlisting[language=c]{sources/codes/b3-securite-a-corriger.c}

    Repérer la ou les erreurs présentes dans ce programme en les \textbf{numérotant
    directement sur le listing} (1, 2\ldots). Puis, dans le cadre
    ci-dessous, donner pour chaque numéro la \textbf{ligne corrigée} et le comportement
    fautif que la correction supprime.

    \evaluationProfGrille{3}{7}
  \end{question}
}

% ==================== PARTIE C — BOUCLES ====================
% Étape de conception : attente de verrouillage, série de sessions, recharge simulée, paliers.
%
% Progression :
%   C1  QCM — sémantique exacte de la borne d'arrêt
%   C2  corriger une boucle (compteur non incrémenté)
%   C3  trace d'une boucle de recharge à gain dégressif (piège : bornes >= atteintes pile à 60 et 80, arrêt pile à 90)
%   C4a compléter l'ossature d'une boucle (3 trous)
%   C4b prédire l'affichage dans un cas limite (0 tour)
%   C5  écrire le corps d'un main (bilan, maximum)

\element{boucle-qcm-borne}{
    \begin{question}{boucle-qcm-borne}
        Au branchement, la borne revérifie le verrouillage du câble toutes les 15 secondes
        tant que la durée d'attente écoulée ne dépasse pas 60 secondes.

        \begin{center}
            \begin{minipage}{.72\textwidth}
                \lstinputlisting[language=c]{sources/codes/c1-qcm-borne.c}
            \end{minipage}
        \end{center}

        Qu'affiche ce fragment ?
        \begin{multicols}{6}
            \begin{reponses}
                \bonne{\texttt{5 75}}
                \mauvaise{\texttt{4 60}}
                \mauvaise{\texttt{4 75}}
                \mauvaise{\texttt{5 60}}
                \mauvaise{\texttt{3 45}}
                \mauvaise{\texttt{4 45}}
            \end{reponses}
        \end{multicols}
    \end{question}
}

\element{boucle-correction}{
    \begin{question}{boucle-correction}
        Pour le bilan d'un point de charge, un développeur inexpérimenté a écrit le programme suivant, qui
        cumule l'énergie délivrée par chacune des sessions de la journée.

        \lstinputlisting[language=c]{sources/codes/c1b-moyenne-a-corriger.c}

        Ce programme ne se comporte pas comme prévu. Décrire le comportement observé à
        l'exécution, puis corriger le programme \textbf{directement sur le code}. Justifier
        ci-dessous.

        \evaluationProfGrille{2}{4}
    \end{question}
}

\element{boucle-trace}{
    \needspace{18\baselineskip}% évite de séparer l'intro du code
    Pour préserver la batterie, la borne réduit la puissance délivrée lorsque le niveau
    de charge augmente : le niveau gagne 20 points par pas de 10 minutes sous 60 \%,
    puis 10 points de 60 à 80 \%, puis 5 points au-delà. Le fragment ci-dessous simule
    cette recharge (\lstinline[language=C]|gain_pct| est le nombre de points de
    pourcentage gagnés pendant un pas de 10 minutes).
    
    \begin{question}{boucle-trace}
        Compléter le tableau en donnant, à la fin de \textbf{chaque tour de boucle}, les
        valeurs de \lstinline[language=C]|gain_pct|, \lstinline[language=C]|niveau_pct| et
        \lstinline[language=C]|minutes|. Numéroter vous-même les tours suivants et
        \textbf{s'arrêter au dernier tour effectivement exécuté} (les lignes en trop restent vides).
        \evaluationProf[2]

        \begin{minipage}[t]{.48\textwidth}
            \lstinputlisting[language=c]{sources/codes/c2-recharge-trace.c}
        \end{minipage}\hfill
        \begin{minipage}[t]{.48\textwidth}~\\
            {\renewcommand{\arraystretch}{2}
            \begin{tabular}{|c|p{1.95cm}|p{1.75cm}|p{1.35cm}|}
                \hline
                \textbf{Tour} & \lstinline[language=C]|gain_pct| & \lstinline[language=C]|niveau_pct| & \lstinline[language=C]|minutes| \\ \hline
                1 & & & \\ \hline
                2 & & & \\ \hline
                  & & & \\ \hline
                  & & & \\ \hline
                  & & & \\ \hline
                  & & & \\ \hline
                  & & & \\ \hline
            \end{tabular}}
        \end{minipage}
    \end{question}

    \begin{question}{boucle-trace-affichage}
        Qu'affiche ce programme après exécution complète de la boucle? 

        \evaluationProfGrille{2}{1}

    \end{question}
}

\element{boucle-completion}{
    \needspace{34\baselineskip}% évite de séparer l'intro du code
    \paragraph{Simulation de la montée en charge : }
    La borne monte le niveau de charge par paliers de 7 points de pourcentage, en partant
    de 15 \%. Le cahier des charges impose que :
    \begin{itemize}\setlength{\itemsep}{0pt}
        \item la boucle s'arrête dès que le niveau atteint \textbf{ou dépasse} le niveau
        cible saisi par l'utilisateur ;
        \item le niveau affiché ne dépasse \textbf{jamais} 100 \%.
    \end{itemize}

    \lstinputlisting[language=c,escapeinside={(*}{*)}]{sources/codes/c3-paliers-a-completer.c}
    
    \begin{question}{boucle-completion}
        Compléter les trois éléments manquants \textbf{directement dans les cadres pointillés}.

        \evaluationProf[2]
    \end{question}
    \begin{question}{boucle-completion-affichage}
        Indiquer ce que le programme affiche si l'utilisateur saisit un niveau cible de 10 \%.
    
        \evaluationProfGrille{2}{2}
    \end{question}
}


\element{boucle-conception}{
    \begin{question}{boucle-conception}
        L'exploitant veut un bilan de fin de journée pour un point de charge.

        Compléter le corps de la fonction \lstinline[language=C]|main| ci-dessous pour qu'elle :
        \begin{itemize}\setlength{\itemsep}{0pt}
            \item demande le nombre de sessions de la journée (entier) ;
            \item demande ensuite, une par une, l'énergie délivrée par chaque session (en kWh,
                  nombre à virgule) ;
            \item calcule l'énergie totale délivrée, le nombre de sessions ayant dépassé
                  20 kWh, et l'énergie de la session la plus importante de la journée ;
            \item affiche ces trois résultats à la fin.
        \end{itemize}

        \textbf{Contraintes :} les valeurs saisies ne doivent pas être conservées toutes en
        mémoire (on ne dispose d'aucune structure permettant de stocker une série de valeurs) ;
        chaque valeur est donc traitée au moment de sa saisie. Le choix de la boucle et la
        gestion correcte du maximum lors de la première itération font partie de l'évaluation.

        \lstinputlisting[language=c]{sources/codes/squelette-main-debut.c}
        \vspace*{-.5em}
        \evaluationProfGrille{4}{12}
        \lstinputlisting[language=c]{sources/codes/squelette-main-fin.c}
    \end{question}
}

% ==================== PARTIE D — FONCTIONS ====================
% Étape de conception : modulariser le logiciel de la borne.
%
% Progression :
%   D1  QCM — portée d'une variable locale (masquage, sans globale)
%   D2  corriger le programme (passage par valeur + retour)
%   D3a QCM multiple — appels acceptés par le compilateur
%   D3b analyse : bonne version et effet des deux autres
%   D4  écrire le corps d'une fonction (décision, codes de retour)

\element{fonc-qcm-portee}{
    \needspace{26\baselineskip}% évite de séparer l'intro du code
    Un développeur tente de compter les sessions de la borne avec le programme suivant :
    
    \begin{center}
        \begin{minipage}{.76\textwidth}
            \lstinputlisting[language=c]{sources/codes/d1-portee.c}
        \end{minipage}
    \end{center}
    
    \begin{question}{fonc-qcm-portee}
        Qu'affiche ce programme ?
        \begin{multicols}{4}
            \begin{reponses}
                \bonne{\texttt{1 1 1 | 0}}
                \mauvaise{\texttt{1 2 3 | 3}}
                \mauvaise{\texttt{1 1 1 | 3}}
                \mauvaise{\texttt{1 2 3 | 0}}
            \end{reponses}
        \end{multicols}
    \end{question}
}

\element{fonc-correction-portee}{
    \begin{question}{fonc-correction-portee}
        Quel problème constatez-vous sur le code de la question précédente ?

        \evaluationProfGrille{3}{6}
    \end{question}
}

\element{fonc-lecture-appel}{
    \needspace{30\baselineskip}% garde code et QCM sur la même page
    Le calcul du montant d'une session est confié à une fonction dont le prototype est :
    
    \begin{center}
        \lstinline[language=C]|float cout_session(float puissance_kw, float duree_h, float tarif_kwh);|
    \end{center}
    
    Quatre appels à cette fonction apparaissent dans le programme :
    
    \begin{center}
        \begin{minipage}{.84\textwidth}
            \lstinputlisting[language=c]{sources/codes/d2-appels-cout.c}
        \end{minipage}
    \end{center}
    
    \begin{questionmult}{fonc-lecture-appel}
        \AMCnoCompleteMulti
        Cocher \textbf{tous} les appels acceptés par le compilateur.
        \begin{multicols}{4}
            \begin{reponses}[o]
                \bonne{appel 1}
                \bonne{appel 2}
                \mauvaise{appel 3}
                \bonne{appel 4}
            \end{reponses}
        \end{multicols}
    \end{questionmult}

    \begin{question}{fonc-lecture-appel-analyse}
        Parmi les appels de la question précédent, indiquer
        lequel est la \textbf{bonne version} pour calculer le coût d'une session de 22 kW pendant 1,5 h à 0,22 €/kWh ?
        \begin{multicols}{4}
            \begin{reponses}[o]
                \bonne{appel 1}
                \mauvaise{appel 2}
                \mauvaise{appel 3}
                \mauvaise{appel 4}
            \end{reponses}
        \end{multicols}
    \end{question}
}

\element{fonc-ecrire-mode}{
    \begin{question}{fonc-ecrire-mode}
        On modularise maintenant la décision d'engagement de la charge.

        On souhaite la fonction
        \lstinline[language=C]|int determiner_mode(float puissance_kw, int cable_verrouille)|
        qui retourne un \textbf{code entier} caractérisant la décision de la borne :
        \begin{itemize}\setlength{\itemsep}{0pt}
            \item si le câble n'est pas verrouillé (\lstinline[language=C]|cable_verrouille| vaut 0),
                  retourner \lstinline[language=C]|-1| quelle que soit la puissance ;
            \item sinon, retourner \lstinline[language=C]|0| si la puissance demandée ne dépasse
                  pas 22 kW, \lstinline[language=C]|1| si elle est comprise entre 22 kW (exclu) et
                  50 kW (inclus), \lstinline[language=C]|2| si elle dépasse 50 kW sans dépasser
                  150 kW ;
            \item si la puissance demandée dépasse 150 kW, retourner \lstinline[language=C]|3|
                  (demande hors domaine de la borne).
        \end{itemize}

        La fonction ne doit rien afficher et doit retourner une valeur dans \textbf{tous} les cas.

        Compléter le corps de la fonction ci-dessous.

        \lstinputlisting[language=c]{sources/codes/d4-mode-debut.c}
        \vspace*{-.5em}
        \evaluationProfGrille{4}{11}
        \lstinputlisting[language=c]{sources/codes/d4-fin.c}
    \end{question}
}

% ==================== PARTIE E — SYNTHÈSE ====================
% Étape de conception : intégration (fonctions + boucle + décisions) et recul critique.
%
%   E1  compléter un programme utilisant determiner_mode (prototype, appel, codes)
%   E2  écrire tarif_applicable (réutilisation possible de determiner_mode)
%   E3  critique de conception

\element{synthese-utilisation}{
    \begin{question}{synthese-utilisation}
        On dispose de la fonction
        \lstinline[language=C]|int determiner_mode(float puissance_kw, int cable_verrouille)|,
        que l'on considère \textbf{correcte}. Elle retourne \lstinline[language=C]|-1| si le
        câble n'est pas verrouillé, \lstinline[language=C]|0| en mode lent,
        \lstinline[language=C]|1| en mode rapide, \lstinline[language=C]|2| en mode
        ultra-rapide et \lstinline[language=C]|3| si la demande dépasse 150 kW. Sa définition
        est placée \textbf{après} le \lstinline[language=C]|main| ; son prototype est rappelé
        en tête du fichier.

        Compléter les cadres pointillés pour que le programme traite les
        \lstinline[language=C]|nb_sessions| demandes de la journée, compte les sessions
        refusées (câble non verrouillé ou demande hors domaine) et affiche le mode des autres.
        Le choix selon le code retourné est réalisé avec une structure
        \lstinline[language=C]|switch|.

        \lstinputlisting[language=c,escapeinside={(*}{*)}]{sources/codes/e1-utilisation-mode-a-completer.c}

        \evaluationProf[3]
    \end{question}
}

\element{synthese-evolution}{
    \begin{question}{synthese-evolution}
        L'exploitant met en place une tarification heures creuses : de 22 h à 5 h inclus, le
        tarif est de 0,15 €/kWh ; le reste du temps, il vaut 0,22 €/kWh jusqu'à 50 kW inclus
        et 0,30 €/kWh au-delà de 50 kW.

        Compléter le corps de la fonction \lstinline[language=C]|tarif_applicable| qui retourne
        le tarif en €/kWh correspondant à une heure de départ de session (entier de 0 à 23) et
        à une puissance donnée. Elle peut utiliser les fonctions déjà écrites dans le sujet,
        par exemple \lstinline[language=C]|determiner_mode| (codes rappelés à la question
        précédente). Elle ne doit rien afficher.

        \lstinputlisting[language=c]{sources/codes/e2-tarif-debut.c}
        \vspace*{-.5em}
        \evaluationProfGrille{4}{10}
        \lstinputlisting[language=c]{sources/codes/d4-fin.c}
    \end{question}
}

\element{synthese-critique}{
    \begin{question}{synthese-critique}
        En reprenant le programme de traitement des sessions de cette partie (utilisation de
        \lstinline[language=C]|determiner_mode|), identifier et expliquer
        \textbf{deux} faiblesses de conception (lisibilité, valeurs numériques écrites en dur,
        robustesse face à des saisies inattendues, absence de vérification, découpage en
        fonctions, comparaison de nombres à virgule\ldots), en indiquant pour chacune comment
        y remédier.

        \evaluationProfGrille{2}{6}
    \end{question}
}

% ==================== PARTIE F — QUESTIONS SUPPLÉMENTAIRES ====================
% Questions plus difficiles déplacées en fin de sujet (retours du référent) :
%   F1  QCM — précision flottante (ex-A2)
%   F2  QCM — évaluation court-circuit de || (ex-B1)

\element{var-qcm-flottant}{
    \begin{question}{var-qcm-flottant}
        La borne cumule trois tranches de facturation de 0,1 € chacune, puis vérifie le total.

        \begin{center}
            \begin{minipage}{.72\textwidth}
                \lstinputlisting[language=c]{sources/codes/a2-qcm-flottant.c}
            \end{minipage}
        \end{center}

        Qu'affiche ce fragment, et pour quelle raison ?
        \begin{reponses}
            \bonne{\texttt{ECART} : 0.1 n'a pas de représentation binaire exacte, le cumul des arrondis rend le total différent de 0.3}
            \mauvaise{\texttt{EXACT} : trois additions de 0.1 donnent exactement 0.3}
            \mauvaise{\texttt{ECART} : un \lstinline[language=C]|float| ne conserve que deux décimales}
            \mauvaise{Rien : comparer deux \lstinline[language=C]|float| avec \lstinline[language=C]|==| provoque une erreur de compilation}
        \end{reponses}
    \end{question}
}

\element{cond-qcm-courtcircuit}{
    \begin{question}{cond-qcm-courtcircuit}
        La fonction de lecture du capteur de tension est lente : chaque lecture est tracée
        à l'écran. Le module de sécurité est écrit ainsi :

        \begin{center}
            \begin{minipage}{.80\textwidth}
                \lstinputlisting[language=c]{sources/codes/b1-court-circuit.c}
            \end{minipage}
        \end{center}

        Qu'affiche exactement ce programme ?
        \begin{reponses}
            \bonne{\texttt{Coupure} seul : l'opérateur \lstinline[language=C]!||! n'évalue pas son second opérande lorsque le premier est vrai}
            \mauvaise{\texttt{[lecture capteur]} puis \texttt{Coupure} : les deux opérandes d'une condition sont toujours évalués}
            \mauvaise{\texttt{[lecture capteur]} seul : la condition est fausse car la tension dépasse la plage autorisée}
            \mauvaise{\texttt{Coupure} seul : un \lstinline[language=C]|printf| placé dans une fonction appelée depuis un \lstinline[language=C]|if| n'est jamais exécuté}
        \end{reponses}
    \end{question}
}


\def\AMCformQuestion#1{{\sc Question #1 :}}

\AMCrandomseed{1237893}
\sujetA
\end{document}
